\documentclass[11pt,a4paper]{article}

% --- Encoding and fonts ---
\usepackage[utf8]{inputenc}
\usepackage[T1]{fontenc}

% --- Mathematics ---
\usepackage{amsmath,amssymb,amsthm}
\usepackage{mathtools}

% --- Page layout ---
\usepackage[margin=1in,headheight=14pt]{geometry}

% --- Hyperlinks ---
\usepackage[colorlinks=true,linkcolor=red,citecolor=blue,urlcolor=blue]{hyperref}

% --- Lists ---
\usepackage{enumitem}

% --- Colors and boxes ---
\usepackage{xcolor}
\usepackage[theorems,skins,breakable]{tcolorbox}

% --- Header/footer ---
\usepackage{fancyhdr}
\pagestyle{fancy}
\fancyhf{}
\fancyhead[L]{\textsc{Lesson 9: Convex Optimization and Duality}}
\fancyhead[R]{\thepage}

% --- Tables ---
\usepackage{booktabs}

% ============================================================
% Theorem environments
% ============================================================
\theoremstyle{definition}
\newtheorem{definition}{Definition}[section]
\newtheorem{example}{Example}[section]
\newtheorem{exercise}{Exercise}[section]

\theoremstyle{plain}
\newtheorem{theorem}{Theorem}[section]
\newtheorem{lemma}[theorem]{Lemma}
\newtheorem{proposition}[theorem]{Proposition}
\newtheorem{corollary}[theorem]{Corollary}

\theoremstyle{remark}
\newtheorem{remark}{Remark}[section]

% ============================================================
% Colored boxes
% ============================================================
\newtcolorbox{keyresult}[1][]{
  colback=green!5!white,
  colframe=green!60!black,
  fonttitle=\bfseries,
  title={Key Result},
  breakable,
  #1
}

\newtcolorbox{rlconnection}[1][]{
  colback=blue!5!white,
  colframe=blue!60!black,
  fonttitle=\bfseries,
  title={RL Connection},
  breakable,
  #1
}

\newtcolorbox{warning}[1][]{
  colback=red!5!white,
  colframe=red!60!black,
  fonttitle=\bfseries,
  title={Warning},
  breakable,
  #1
}

% ============================================================
% Custom commands
% ============================================================
\newcommand{\R}{\mathbb{R}}
\newcommand{\N}{\mathbb{N}}
\newcommand{\Q}{\mathbb{Q}}
\newcommand{\Z}{\mathbb{Z}}
\newcommand{\C}{\mathbb{C}}
\newcommand{\Prob}{\mathbb{P}}
\newcommand{\E}{\mathbb{E}}
\newcommand{\indicator}{\mathbf{1}}

\DeclareMathOperator{\dom}{dom}
\DeclareMathOperator{\epi}{epi}
\DeclareMathOperator{\conv}{conv}
\DeclareMathOperator{\prox}{prox}
\DeclareMathOperator*{\argmin}{arg\,min}
\DeclareMathOperator*{\argmax}{arg\,max}
\DeclareMathOperator{\tr}{tr}
\DeclareMathOperator{\diag}{diag}
\DeclareMathOperator{\interior}{int}

% ============================================================
\title{Lesson 25: Convex Optimization Problems}
\author{%
  Session 9 of 102 \quad$\mid$\quad Phase I: Mathematical Foundations \quad$\mid$\quad 8h \\[4pt]
  \textit{Primary texts: Boyd \& Vandenberghe Ch.\,4--5, 9--11; Nesterov Ch.\,1--4}
}
\date{}

\begin{document}
\maketitle
\tableofcontents
\newpage

\setcounter{section}{0}

\section{Introduction}\label{sec:intro}

Optimization is the computational engine beneath virtually every algorithm in
machine learning and reinforcement learning. Policy gradient methods solve
surrogate optimization problems at each iteration; value function approximation
minimizes projected Bellman error; and constrained MDPs reduce to saddle-point
problems whose structure is governed by duality theory.

This document develops the theory of convex optimization from first principles:
the hierarchy of convex programs (Section~\ref{sec:convex-problems}), Lagrangian
duality and the passage from primal to dual (Section~\ref{sec:duality}), the
Karush--Kuhn--Tucker conditions (Section~\ref{sec:kkt}), gradient descent and
its convergence guarantees (Section~\ref{sec:gd}), Nesterov's accelerated method
(Section~\ref{sec:nesterov}), proximal methods for non-smooth objectives
(Section~\ref{sec:proximal}), and interior-point methods
(Section~\ref{sec:interior}). We close with exercises
(Section~\ref{sec:exercises}) and a summary table of convergence rates
(Section~\ref{sec:summary}).

Throughout, we assume familiarity with real analysis (continuity, compactness)
and linear algebra (eigenvalues, positive definiteness) at the level of the
Phase~00 prerequisites.

%% ============================================================
%% SECTION 2: Convex Optimization Problems
%% ============================================================

\section{Convex Optimization Problems}\label{sec:convex-problems}

\subsection{Standard Form}

\begin{definition}[Convex optimization problem]\label{def:convex-opt}
A \emph{convex optimization problem} in standard form is
\begin{equation}\label{eq:std-form}
\begin{aligned}
  \min_{x \in \R^n} \quad & f_0(x) \\
  \text{subject to} \quad & f_i(x) \le 0, \quad i = 1,\dots,m, \\
                          & a_j^\top x = b_j, \quad j = 1,\dots,p,
\end{aligned}
\end{equation}
where $f_0, f_1, \dots, f_m : \R^n \to \R$ are convex functions and the
equality constraints are affine.
\end{definition}

\begin{definition}[Feasible set, optimal value]\label{def:feasible}
The \emph{feasible set} is
$\mathcal{F} = \{x \in \R^n : f_i(x) \le 0,\; a_j^\top x = b_j\}$.
The \emph{optimal value} is $p^* = \inf\{f_0(x) : x \in \mathcal{F}\}$.
A point $x^*$ with $f_0(x^*) = p^*$ is an \emph{optimal point} (or
\emph{minimizer}).
\end{definition}

\begin{theorem}[Local optimality implies global optimality]\label{thm:local-global}
For the convex program~\eqref{eq:std-form}, every local minimizer is a global
minimizer.
\end{theorem}

\begin{proof}
Let $x^*$ be a local minimizer: there exists $\delta > 0$ such that
$f_0(x^*) \le f_0(x)$ for all feasible $x$ with $\|x - x^*\| \le \delta$.
Suppose for contradiction there exists a feasible $y$ with $f_0(y) < f_0(x^*)$.
Consider $z_t = (1-t)x^* + ty$ for $t \in (0,1)$. By convexity of each $f_i$
and affinity of the equality constraints, $z_t$ is feasible. Moreover,
\[
  f_0(z_t) \le (1-t) f_0(x^*) + t f_0(y) < f_0(x^*)
\]
for all $t \in (0,1)$. Choosing $t$ small enough that
$\|z_t - x^*\| = t\|y - x^*\| < \delta$ contradicts local optimality.
\end{proof}

\subsection{Important Special Cases}

\begin{definition}[Linear program -- LP]\label{def:lp}
A \emph{linear program} has the form
\[
  \min_{x} \; c^\top x \quad \text{s.t.} \quad Ax \le b, \; Cx = d,
\]
where all functions are affine (hence convex).
\end{definition}

\begin{definition}[Quadratic program -- QP]\label{def:qp}
A \emph{quadratic program} minimizes a convex quadratic over a polyhedral
feasible set:
\[
  \min_{x} \; \tfrac{1}{2} x^\top P x + q^\top x \quad \text{s.t.} \quad Ax \le b,
\]
with $P \succeq 0$ (positive semidefinite).
\end{definition}

\begin{definition}[Second-order cone program -- SOCP]\label{def:socp}
An SOCP has constraints of the form $\|A_i x + b_i\|_2 \le c_i^\top x + d_i$
(second-order cone constraints), which are convex.
\end{definition}

\begin{definition}[Semidefinite program -- SDP]\label{def:sdp}
An SDP minimizes a linear objective subject to a linear matrix inequality:
\[
  \min_{x} \; c^\top x \quad \text{s.t.} \quad F_0 + \sum_{i=1}^n x_i F_i \succeq 0,
\]
where $F_0, \dots, F_n$ are symmetric matrices.
\end{definition}

\begin{remark}[Hierarchy]\label{rem:hierarchy}
We have the strict inclusion chain:
$\text{LP} \subset \text{QP} \subset \text{SOCP} \subset \text{SDP} \subset
\text{convex programs}$.
\end{remark}

\begin{rlconnection}
In RL, linear programs arise in the LP formulation of MDPs (minimize
$\sum_s c(s) V(s)$ subject to Bellman feasibility constraints). Semidefinite
programs appear in robust MDP formulations and certain Lyapunov-based safety
constraints. Quadratic programs arise in trust-region policy optimization when
the surrogate objective is approximated to second order.
\end{rlconnection}

%% ============================================================
%% SECTION 3: Lagrangian Duality
%% ============================================================


\end{document}